% ======================================================================
%  Section 4.2.2 --- The weighted control law.  Corrected.
%
%  THREE SUBSTANTIVE PROBLEMS, then four smaller ones.
%
%  A. Equation (4.4) does not match the implementation, and another section
%     already refers to the missing piece. The code computes
%         H = ( mu*diag(H_D) + H_D + muFloor*I )^{-1}
%     with muFloor = max(1e-6 * mean(diag H_D), 1e-12). Section 4.3.1 ends
%     with "what makes the safeguard recorded there necessary rather than
%     precautionary" -- but no safeguard is recorded here. Section 4.7.1
%     likewise attributes invertibility to "the fixed floor added to the
%     damped Hessian". The term has to appear in (4.4).
%
%  B. The null-space argument is too strong as stated. "the damped matrix
%     inherits the null space of H_D" does not follow in general: a
%     rank-deficient positive semi-definite matrix can have every diagonal
%     entry positive, [[1,1],[1,1]] being the standard example, and then
%     mu*diag(H_D) is positive definite and the damping regularises
%     perfectly well. The statement is true in exactly one case -- when a
%     diagonal entry of H_D vanishes -- and that case is worth stating
%     precisely, because for a positive semi-definite matrix a zero diagonal
%     entry forces the whole row and column to vanish, which is exactly what
%     the weighting does when it annihilates a direction. Corrected below.
%
%  C. The D-squared convention is flagged but its consequence is not drawn.
%     Section 4.2.1 derives the weights of an IRLS problem min sum w_j e_j^2,
%     whose normal equations carry W = diag(w). Equation (4.4) carries D^2 =
%     W^2. So the scheme realises the IRLS problem with weights w_j^2, not
%     w_j. This is the standard construction in visual servoing and is
%     harmless, but it should be said rather than left for a reader in
%     estimation theory to notice.
% ======================================================================

\subsection{The weighted control law}
\label{sec:weighted-law}
%%% Sections 4.7.2 and 4.7.3 refer back here; give the subsection a label.

Let $\mathbf{D}=\mathrm{diag}(w_{1},\dots,w_{P})$ collect the weights, a
symmetric positive semi-definite matrix by construction. Introducing
$\mathbf{D}$ into the Levenberg--Marquardt control law~\eqref{eq:lm-law} of
Chapter~\ref{ch:herpvs} gives the robust control law used throughout this
chapter,
%%% (A) The epsilon term is the safeguard the implementation applies and
%%% that Sections 4.3.1 and 4.7.1 already refer to. Without it the equation
%%% in the thesis is not the equation in the program.
\begin{equation}
    \mathbf{V}_{c}=-\lambda\,\Bigl(\mathbf{H}_{\mathbf{D}}
    +\mu\,\mathrm{diag}\bigl(\mathbf{H}_{\mathbf{D}}\bigr)
    +\varepsilon\,\mathbf{I}_{6}\Bigr)^{-1}
    \mathbf{L}_{R}^{\top}\mathbf{D}^{2}\,\mathbf{e},
    \qquad
    \mathbf{H}_{\mathbf{D}}=\bigl(\mathbf{D}\mathbf{L}_{R}\bigr)^{\top}
    \bigl(\mathbf{D}\mathbf{L}_{R}\bigr)
    =\mathbf{L}_{R}^{\top}\mathbf{D}^{2}\mathbf{L}_{R},
    \label{eq:robust-lm-law}
\end{equation}
in which $\mathbf{L}_{R}$ is the stacked four-scale Hermite interaction
matrix~\eqref{eq:stacked-L}, the gain $\lambda$ and damping $\mu$ retain the
adaptive form~\eqref{eq:lambda}--\eqref{eq:mu}, and
\begin{equation}
    \varepsilon=\max\!\left(10^{-6}\,\frac{\operatorname{tr}\mathbf{H}_{\mathbf{D}}}{6},
    \;10^{-12}\right)
    \label{eq:hessian-floor}
\end{equation}
is a fixed floor on the damped Hessian whose purpose is explained below. The
weighting acts simultaneously on the error and on the interaction matrix,
which is what distinguishes it from a simple masking of the error: a
measurement judged unreliable contributes neither to the descent direction nor
to the curvature estimate.

\paragraph{Why the floor is necessary.}
The inverse in~\eqref{eq:robust-lm-law} requires that the surviving weights
leave $\mathbf{H}_{\mathbf{D}}$ of rank six. This is a stronger requirement
than in the unweighted case, for a reason specific to the form of the damping.
%%% (B) The precise statement. A rank-deficient PSD matrix need NOT have a
%%% zero diagonal entry, so the damping does not in general fail; it fails in
%%% exactly the case described here, which is also exactly the case robust
%%% rejection produces.
Suppose the weighting annihilates some direction of the feature space
entirely, so that the $i$-th diagonal entry of $\mathbf{H}_{\mathbf{D}}$
vanishes. Since $\mathbf{H}_{\mathbf{D}}$ is positive semi-definite, a zero
diagonal entry forces the whole of its $i$-th row and column to vanish with
it; the matrix $\mathrm{diag}(\mathbf{H}_{\mathbf{D}})$ then has a zero in the
same place, so the damping term $\mu\,\mathrm{diag}(\mathbf{H}_{\mathbf{D}})$
vanishes on that direction as well and the damped matrix is singular. Scaling
the damping by $\mathrm{diag}(\mathbf{H}_{\mathbf{D}})$ rather than by the
identity, which is what makes the scheme invariant to the scaling of the
individual degrees of freedom, is precisely what deprives it of the ability to
regularise in this case. A redescending estimator that rejects too large a
fraction of the measurements can therefore render the control law undefined
rather than merely ill-conditioned.

The floor $\varepsilon$ of~\eqref{eq:hessian-floor} removes this failure mode
at negligible cost. It is six orders of magnitude below the mean diagonal of
$\mathbf{H}_{\mathbf{D}}$, so it does not perturb the step while the Hessian
is well conditioned, and it bounds the inverse away from singularity when a
direction is lost. The absolute floor of $10^{-12}$ covers the degenerate case
in which $\mathbf{H}_{\mathbf{D}}$ vanishes altogether.
%%% If you would rather keep the floor out of the main equation, the
%%% alternative is to state it here as an implementation detail and cite it
%%% from 4.3.1 and 4.7.1. I recommend against that: the numbers reported in
%%% 4.7 come from a program that applies it at every iteration, and it is
%%% the reason the rank condition is a theoretical rather than a practical
%%% constraint.

\paragraph{Two conventions.}
The weights are evaluated from the residuals of the current servoing iteration
and held fixed within the step, consistently with the iteratively reweighted
scheme of Section~\ref{sec:least-squares}.
%%% Was sec:robust-framework, which is the whole of 4.2; IRLS is described in
%%% 4.2.1 specifically.
And they enter~\eqref{eq:robust-lm-law} by scaling both the interaction matrix
and the error before the normal equations are formed, so that $\mathbf{D}$
appears squared in $\mathbf{H}_{\mathbf{D}}$ --- the same construction used
in~\cite{key26}. The weight $w_{j}$ of~\eqref{eq:irls-weight} is applied
directly, without a compensating square root, so that the normal equations
carry $w_{j}^{2}$ and the scheme realises the reweighted problem
$\min\sum_{j}w_{j}^{2}e_{j}^{2}$ rather than the
$\min\sum_{j}w_{j}e_{j}^{2}$ of Section~\ref{sec:least-squares}.
%%% (C) This is the sentence that was missing. The consequence is worth
%%% stating and then dismissing, which the next sentence does.
Since $w_{j}\in[0,1]$ and the map $w\mapsto w^{2}$ is monotone and fixes both
endpoints, the sets of rejected and of fully retained measurements are
identical under the two conventions and only the grading between them differs;
the choice follows the visual servoing literature and is recorded here for
reproducibility rather than because anything depends on it.

Two properties of~\eqref{eq:robust-lm-law} are worth stating explicitly, as
they structure the remainder of the chapter.

\begin{remark}[The unweighted scheme is the special case $\mathbf{D}=\mathbf{I}$]
\label{rem:baseline}
Setting all weights to unity reduces $\mathbf{H}_{\mathbf{D}}$ to
$\mathbf{L}_{R}^{\top}\mathbf{L}_{R}$ and~\eqref{eq:robust-lm-law} to the
control law~\eqref{eq:lm-law} of Chapter~\ref{ch:herpvs}, exactly. The
unweighted scheme is therefore not a different method introduced for
comparison, but the degenerate case of the present formulation. This is what
makes the first arm of the ablation of Section~\ref{sec:ablation-protocol} a
genuine baseline: the three variants differ in the definition of $\mathbf{D}$
and in nothing else.
\end{remark}

\begin{remark}[What the weights may depend upon]
\label{rem:weight-information}
The estimator~\eqref{eq:irls-weight} determines $w_{j}$ from the residual
$e_{j}$ alone. Nothing in the weighted control law~\eqref{eq:robust-lm-law}
requires this: any non-negative diagonal $\mathbf{D}$ preserving the rank
condition above yields a valid weighted problem. The weights may therefore be
made to depend on further information about the measurement, and in particular
on the local image structure at the position where it is taken. Exploiting
this freedom is the subject of Section~\ref{sec:hermite-weight}.
\end{remark}

% ======================================================================
%  D. THE REMARK ENVIRONMENT IS NOT RENDERING
%
%  In the PDF these two remarks appear as running text with their titles in
%  literal square brackets:
%      [The unweighted scheme is the special case D = I] Setting all ...
%      [What the weights may depend upon] The estimator ...
%  so the optional argument is being typeset rather than used as a heading.
%  That is also why \ref{rem:baseline} prints as "Remark ??" on p. 75 while
%  resolving on p. 66. Declare the environment properly in the preamble:
%
%      \usepackage{amsthm}
%      \theoremstyle{remark}
%      \newtheorem{remark}{Remark}[section]
%
%  and rebuild twice so the .aux settles. Both remarks are cited by number
%  from 4.4.1, 4.6.1 and 4.7.1, so this is not cosmetic.
%
%  E. SMALLER POINTS
%  - "regularizing" is -ize; the chapter body is -ise elsewhere.
%  - \eqref{eq:lm-law} and \eqref{eq:stacked-L} are Chapter 3 labels;
%    confirm they still resolve, since 4.4.2 currently prints Section ?? for
%    its Chapter 3 cross-references.
%  - The epsilon of (4.5) collides with no existing symbol in the chapter,
%    but check Chapter 3, where epsilon is sometimes a convergence
%    threshold. If it is taken, use \varepsilon_{H} or \delta.
% ======================================================================
